\section{Positioning against the four literature lines}

The paper's questions --- existence, identifiability, realisation form, dynamics, construction --- decompose across four lines of work, and no single line contains the join. This appendix states, line by line, what each supplies, what it does not contain, and where a reader should look for the nearest collision with our claims. It is a table, not an argument; the argument is \S{}2.


\begin{table*}[t]
\centering
\scriptsize
\caption*{\textbf{Table G.1.} Positioning table. "Nearest collision" names the work a reviewer is most likely to raise against the corresponding claim of this paper, with the reason it does not anticipate it.}
\setlength{\tabcolsep}{2pt}
\begin{tabularx}{\textwidth}{@{}>{\raggedright\arraybackslash}X >{\raggedright\arraybackslash}X >{\raggedright\arraybackslash}X >{\raggedright\arraybackslash}X >{\raggedright\arraybackslash}X >{\raggedright\arraybackslash}X >{\raggedright\arraybackslash}X@{}}
\toprule
\textbf{line} & \textbf{representative works} & \textbf{central question in that line} & \textbf{what it supplies to this paper} & \textbf{what it does not contain} & \textbf{our increment} & \textbf{nearest collision} \\
\midrule
\textbf{Equivariance measurement} & Lenc and Vedaldi [1]; Gruver et al. [2]; Bruintjes et al. [3]; Romero and Lohit [4] & how equivariant is a given representation, and can a probe-space map make it more so? & the empirical object --- a frozen representation whose response to a known transformation can be measured --- and the practice of scoring it against a re-rendered reference & no information condition: equivariance is scored, never asked to exist; no account of which part of the algebra the score conflates; no composition behaviour & existence is an information condition on the encoder's fibres (\S{}3.2), priced exactly in the linear layer (Thms 1--2) and with the fixed-linear class's score exact in the rectifier layer (Props C--E); and a family of operators can be perfectly self-consistent while failing the transformation (\S{}5.4) & Gruver et al.'s Lie-derivative score is a \textit{local} equivariance error; it is silent on the global existence question, and our Proposition D shows the local quantity it measures belongs to the gate-change term, whose role in the failure is decided by the region-transition criterion \\
\textbf{Symmetry-based and equivariant architectures} & Cohen and Welling [9]; Weiler and Cesa [10]; Lengyel et al. [11]; Yang et al. [12]; Yang et al. [13]; Higgins et al. [15] & how do we \textit{build} a representation with a prescribed symmetry, and what can such a network express? & the representation-theoretic candidate forms --- invariant blocks and harmonic rotation blocks (Prop, \S{}3.6) --- and the demonstration that exact architectural equivariance is achievable when the symmetry is known a priori & the transformation must be known and built in; nothing about the structure a \textit{pretrained, opaque} network already carries; the group-versus-semigroup distinction is assumed rather than measured & the measured structure \textit{is} the blueprint: the paper reads the organisation off frozen networks and then builds the interface from the measurement (\S{}4, \S{}7), with the action never fitted & CEConv and its successors assume colour equivariance is desirable and impose it; we show that on the colour instance it is \textit{already approximately present} in ordinary backbones, and measure where it is not \\
\textbf{Colour and perceptual transformations} & Koenderink [17]; Duits et al. [18]; colour-appearance literature & how should a colour or scale transformation be represented analytically so that it is well behaved? & the reference transformations themselves --- the hue rotation and the diffusion semigroup --- and their algebraic character as implemented physics & no network-side condition; nothing about whether a learned representation carries the transformation, composes it, or can be edited through it & the same transformations used as ground-truth probes on frozen networks, with the algebra (group vs semigroup) as the \textit{independent variable} of the measurement & the dissipative family is the sharp case: the truth action composes to $1.8\times10^{-7}$, yet no basis tested carries a shared decay rate, so the carrier fails for a reason the analytic theory alone does not predict (L15) \\
\textbf{Causal abstraction and latent linearisation} & Geiger et al. [19], [20]; Park et al. [21]; Nadaf [22]; Venkatesh and Kurapath [23] & when do high-level interventions correspond to low-level ones, and when is a latent direction a well-defined intervention? & the interventional-consistency frame, of which our fibre condition is the transformation-action specialisation, and the caution that latent directions are not identifiable without a stated reference & single interventions rather than a whole algebra; no composition law and no price for violating it; no realisation-form question; no construction & the \textit{whole} algebra must descend (\S{}3.2), composition is priced quantitatively (Thm 4), identifiability is carried by the kernel and stabilisers of the representation (\S{}3.1), and the question continues to family capacity (Thm 3) and construction (\S{}7) & our fibre condition is formally Geiger et al.'s consistency condition specialised to transformation actions --- we say so in \S{}2.4 and claim only the specialisation's additional content: the algebra, the pricing, and the construction \\
\textbf{Dynamical / Koopman-style linearisation} & Koopman-style embedding work as surveyed in \S{}2.4 & when does a nonlinear system admit an invariant finite-dimensional subspace on which the evolution is linear? & the idea that a linear operator on an embedding can stand for a nonlinear action, which our band-limited global family instantiates in truncated form & the subspace is \textit{designed or learned} and the dynamics are temporal; the line does not ask whether a \textit{given, frozen} representation admits a linear realisation, nor what obstructs one & the realisation question is \textit{decided} for the given representation, not assumed: existence is an iff condition (Thm 1), the closure criterion decides which carrier coordinates can carry a fixed action (Props F--G), the capacity of the natural per-dimension class is bounded before fitting (Thm 3), and the failure is attributed to a measurable term (Prop D) & a truncated Koopman operator on features is one of our rejected global families (\S{}5.2) unless the orbit is band-limited --- the measured spectrum ($k\le2$ at $84$--$88\%$) is what makes the truncated form work \\
\bottomrule
\end{tabularx}
\end{table*}

\subsubsection*{G.1 Where this paper does \textit{not} compete}

Three boundaries are worth stating so that the positioning is not read as a claim of priority.

\setcounter{enumi}{0}
\begin{enumerate}
\item \textbf{Architecture design.} The paper does not propose a new layer, loss, or training scheme, and no claim here depends on one. The interface of \S{}7 is a read-out and a fixed action on \textit{existing} features; where it is trained (the read-in $E_\theta$), it is trained under ordinary supervision, and the paper's contribution is that the \textit{form} of the action is fixed by measurement rather than learned.
\item \textbf{Perceptual theory.} The hue rotation and the diffusion semigroup are used as exact, re-renderable references (Assumption A1); the paper takes their physical validity in the rendered domain as given, and treats questions of colour appearance, spectral rendering, or perceptual uniformity as out of scope (\S{}2.3).
\item \textbf{A universal claim about all networks.} Every general statement here is a statement about encoders and algebras under stated assumptions (Appendix B); every \textit{quantitative} statement is a measurement on the site and family it names. The claim--scope table (\S{}3.7) and Appendix A fix this per row, and \S{}8.3 states the two places where the theory's magnitude does not carry over.
\end{enumerate}

